% SOURCE_AUTHORITY: sga5_fr_workpass.tex, lines 7578--7632 (LF-normalized unit).
% SOURCE_SHA256: 791F4EFFC5E02832D5D77ED03518C8156D6F07E4C8238B03545DB93D883FBB28
\subsection*{5. Sistemas proyectivos \(J\)-ádicos y AR-\(J\)-ádicos noetherianos}

En todo este apartado se dan un anillo conmutativo noetheriano \(A\), un ideal \(J\) de \(A\) y una \(A\)-categoría abeliana \(\mathcal C\).

\subsubsection*{5.1. Generalidades}

\begin{statement}{Definición 5.1.1}
Un objeto \(X\) de la categoría \(\Hom(\mathbb N^\circ,\mathcal C)\) se dice \(J\)-ádico noetheriano (resp. artiniano) si es \(J\)-ádico y sus componentes son objetos noetherianos (resp. artinianos) de la categoría \(\mathcal C\).
\end{statement}

Se denotará por \(J\)-adn\((\mathcal C)\) (resp. \(J\)-adf\((\mathcal C)\)) la subcategoría plena de
\(\mathcal P=\Hom(\mathbb N^\circ,\mathcal C)\) cuyos objetos son los sistemas proyectivos \(J\)-ádicos noetherianos (resp. artinianos). Es claro (3.1.3) que una extensión de sistemas proyectivos \(J\)-ádicos noetherianos (resp. artinianos) es del mismo tipo.

\begin{statement}{Proposición 5.1.2}
Sea \(X\) un objeto \(J\)-ádico de \(\mathcal P\). Para que \(X\) sea \(J\)-ádico noetheriano (resp. artiniano), es necesario y suficiente que sea \(J\)-ádico y que \(X_0\) sea noetheriano (resp. artiniano).
\end{statement}

\begin{prooftext}
Sea, en efecto, \(X\) un sistema proyectivo \(J\)-ádico tal que \(X_0\) sea noetheriano (resp. artiniano). Consideremos entonces el epimorfismo canónico (4.2.4):
\[
\gr_J(A)\otimes_A X_0\longrightarrow \grs(X).
\]
Como \(X_0\) es noetheriano (resp. artiniano) y \(J^n\) es de tipo finito para todo \(n\), la existencia de este epimorfismo muestra que, para todo entero \(n\), el núcleo del morfismo canónico
\[
X_n\longrightarrow X_{n-1}
\]
es noetheriano (resp. artiniano). Se deduce de inmediato por inducción sobre \(n\) que los \(X_n\) son noetherianos (resp. artinianos).
\end{prooftext}

\begin{statement}{Definición 5.1.3}
Sea \(X\) un sistema proyectivo indexado por \(\mathbb N\) con valores en \(\mathcal C\). Se dice que es AR-\(J\)-ádico noetheriano (resp. artiniano) si es AR-\(J\)-ádico y sus componentes son objetos noetherianos (resp. artinianos).
\end{statement}

Se denotará por AR-\(J\)-adn\((\mathcal C)\) (resp. AR-\(J\)-adf\((\mathcal C)\)) la subcategoría plena de \(\Hom_{\mathrm{AR}}(\mathbb N^\circ,\mathcal C)\) (2.4) engendrada por las imágenes de los sistemas proyectivos AR-\(J\)-ádicos noetherianos (resp. artinianos).

Entonces es claro que los funtores canónicos
\[
P^n_{\mathrm{AR}}:J\text{-adn}(\mathcal C)\longrightarrow \mathrm{AR}\text{-}J\text{-adn}(\mathcal C),
\]
\[
P^f_{\mathrm{AR}}:J\text{-adf}(\mathcal C)\longrightarrow \mathrm{AR}\text{-}J\text{-adf}(\mathcal C),
\]
obtenidos por restricción de \(P_{\mathrm{AR}}\) (2.4), son equivalencias de categorías.

Recordemos brevemente el enunciado y la demostración del teorema de Hilbert.

\begin{statement}{Teorema 5.1.4}
Teorema de Hilbert. Sean \(A\) un anillo conmutativo, \(\mathcal C\) una \(A\)-categoría abeliana y \(M\) un objeto de tipo finito (resp. noetheriano) de \(\mathcal C\). Para toda \(A\)-álgebra de tipo finito \(B\), graduada por \(\mathbb N\), el \((B,\mathcal C)\)-módulo graduado
\[
B\otimes_A M
\]
es de tipo finito (resp. noetheriano).

N.B. Si \(B=\bigoplus_{n\ge0}B_n\), \(B\otimes_A M\) designa el objeto graduado \(\bigl(B_n\otimes_A M\bigr)_{n\in\mathbb N}\).
\end{statement}
